% Options for packages loaded elsewhere
\PassOptionsToPackage{unicode}{hyperref}
\PassOptionsToPackage{hyphens}{url}
%
\documentclass[
  11pt,
]{article}
\usepackage{amsmath,amssymb}
\usepackage{iftex}
\ifPDFTeX
  \usepackage[T1]{fontenc}
  \usepackage[utf8]{inputenc}
  \usepackage{textcomp} % provide euro and other symbols
\else % if luatex or xetex
  \usepackage{unicode-math} % this also loads fontspec
  \defaultfontfeatures{Scale=MatchLowercase}
  \defaultfontfeatures[\rmfamily]{Ligatures=TeX,Scale=1}
\fi
% \usepackage{lmodern}
\ifPDFTeX\else
  % xetex/luatex font selection
\fi
% Use upquote if available, for straight quotes in verbatim environments
\IfFileExists{upquote.sty}{\usepackage{upquote}}{}
\IfFileExists{microtype.sty}{% use microtype if available
  % \usepackage[]{microtype}
  % \UseMicrotypeSet[protrusion]{basicmath} % disable protrusion for tt fonts
}{}
\makeatletter
\@ifundefined{KOMAClassName}{% if non-KOMA class
  \IfFileExists{parskip.sty}{%
    \usepackage{parskip}
  }{% else
    \setlength{\parindent}{0pt}
    \setlength{\parskip}{6pt plus 2pt minus 1pt}}
}{% if KOMA class
  \KOMAoptions{parskip=half}}
\makeatother
\usepackage{xcolor}
\usepackage{longtable,booktabs,array}
\usepackage{calc} % for calculating minipage widths
% Correct order of tables after \paragraph or \subparagraph
\usepackage{etoolbox}
\makeatletter
\patchcmd\longtable{\par}{\if@noskipsec\mbox{}\fi\par}{}{}
\makeatother
% Allow footnotes in longtable head/foot
\IfFileExists{footnotehyper.sty}{\usepackage{footnotehyper}}{\usepackage{footnote}}
\makesavenoteenv{longtable}
\setlength{\emergencystretch}{3em} % prevent overfull lines
\providecommand{\tightlist}{%
  \setlength{\itemsep}{0pt}\setlength{\parskip}{0pt}}
\setcounter{secnumdepth}{-\maxdimen} % remove section numbering
\usepackage{amsmath,amssymb,amsfonts}
\usepackage{mathtools}
\usepackage{booktabs}
\usepackage{longtable}
\usepackage{array}
\usepackage{float}
\usepackage[margin=1in]{geometry}
\usepackage{hyperref}
\hypersetup{colorlinks=true,linkcolor=blue,citecolor=blue,urlcolor=blue}
\usepackage{parskip}
\setlength{\parskip}{0.5em}
\usepackage{enumitem}
\ifLuaTeX
  \usepackage{selnolig}  % disable illegal ligatures
\fi
\IfFileExists{bookmark.sty}{\usepackage{bookmark}}{\usepackage{hyperref}}
\IfFileExists{xurl.sty}{\usepackage{xurl}}{} % add URL line breaks if available
\urlstyle{same}
\hypersetup{
  hidelinks,
  pdfcreator={LaTeX via pandoc}}

\author{}
\date{}

\begin{document}

\hypertarget{dynamics-selection-and-emergent-general-relativity-in-the-observational-incompleteness-framework}{%
\section{Dynamics Selection and Emergent General Relativity in the
Observational Incompleteness
Framework}\label{dynamics-selection-and-emergent-general-relativity-in-the-observational-incompleteness-framework}}

\hypertarget{alex-maybaum}{%
\subsubsection{Alex Maybaum}\label{alex-maybaum}}

\hypertarget{march-2026}{%
\subsubsection{March 2026}\label{march-2026}}

\begin{center}\rule{0.5\linewidth}{0.5pt}\end{center}

\hypertarget{abstract}{%
\subsection{Abstract}\label{abstract}}

The Observational Incompleteness (OI) framework {[}1{]} proves that
quantum mechanics emerges necessarily from embedded observation in a
deterministic system. This paper tests the framework's rigidity by
constructing an explicit lattice realization and asking: does the
dynamics uniquely selected by the QM requirement also produce the inputs
for general relativity? It does. Among all second-order reversible
nearest-neighbor dynamics on a lattice, center independence (necessary
for emergent QM), spatial isotropy, and linearity uniquely select the
discrete wave equation --- for any alphabet size q \ensuremath{\geq} 2 and dimension d \ensuremath{\geq}
1. This dynamics, without further tuning, produces area-law entropy,
Lorentz-invariant dispersion, horizon thermality, and all inputs to
Jacobson's thermodynamic derivation of Einstein's equations. The
derivation chain has seven links, all analytically proven for both the
linear wave equation over \ensuremath{\mathbb{R}} and the mod-q wave equation over \ensuremath{\mathbb{Z}}/q\ensuremath{\mathbb{Z}}. Each
link is an independent check that could have failed; none does. The
structural consistency --- a dynamics selected by one criterion
(emergent QM) passing seven additional tests (emergent GR) with no free
parameters --- constitutes corroborative evidence for the framework's
central claim.

\begin{center}\rule{0.5\linewidth}{0.5pt}\end{center}

\hypertarget{introduction}{%
\subsection{1. Introduction}\label{introduction}}

The companion paper {[}1{]} proves that any observer embedded in a
deterministic, bijective system with a coupled, slow, high-capacity
hidden sector necessarily describes the visible sector using
P-indivisible stochastic dynamics, equivalent to unitary quantum
mechanics. The characterization theorem establishes that QM and embedded
observation are equivalent --- the conditions are necessary and
sufficient. (The result is established by two independent routes: the
Barandes stochastic-quantum correspondence {[}2{]} and a Stinespring
dilation construction {[}1, Appendix A{]}.)

A natural question follows: is the framework rigid or flexible? A
framework that accommodates any dynamics equally well has little
predictive content. A framework that constrains the dynamics --- and
where the constrained dynamics independently reproduces known physics
beyond QM --- provides structural evidence for its correctness.

This paper addresses that question. We construct a specific system ---
the mod-q discrete wave equation on a lattice with a checkerboard
partition --- and show that it is uniquely selected by the requirement
of emergent QM plus two physical symmetries (isotropy and linearity). We
then check whether this dynamics, without modification, provides the
inputs needed for Jacobson's thermodynamic derivation of Einstein's
equations {[}3{]}. The construction proceeds through seven analytical
links. At each link, the framework could have produced the wrong answer:
the entropy could have scaled with volume rather than area, the
dispersion relation could have been non-relativistic, the horizon state
could have been non-thermal. None of these failures occurs. The wave
equation passes every check, with no free parameters.

The paper's primary contribution is therefore not ``deriving GR'' ---
the thermodynamic derivation is due to Jacobson {[}3{]} --- but
demonstrating that the OI framework is overconstrained: the dynamics
selected by QM emergence uniquely and correctly produces the inputs for
GR emergence, providing independent corroborative evidence for the
framework.

\begin{center}\rule{0.5\linewidth}{0.5pt}\end{center}

\hypertarget{the-system}{%
\subsection{2. The System}\label{the-system}}

The mod-q discrete wave equation on a d-dimensional cubic lattice \ensuremath{\Lambda} =
(\ensuremath{\mathbb{Z}}/L\ensuremath{\mathbb{Z}})\^{}d with periodic boundary conditions:

\[x_\mathbf{i}(t+1) = \left(\sum_{\mathbf{j} \sim \mathbf{i}} x_\mathbf{j}(t) \ - \ x_\mathbf{i}(t-1)\right) \bmod q\]

where the sum runs over the 2d nearest neighbors of site \textbf{i}, and
q is a positive integer.

\textbf{Axiom 1 (Bijectivity).} The phase-space map F: (x(t\ensuremath{-}1), x(t)) \ensuremath{\to}
(x(t), x(t+1)) is bijective.

\textbf{Axiom 2 (Finiteness).} The configuration space C = \{0, 1,
\ldots, q\ensuremath{-}1\}\^{}\{\textbar \ensuremath{\Lambda}\textbar\} is finite, with
\textbar C\textbar{} = q\textsuperscript{\{L}d\}.

\textbf{Axiom 3 (Partition).} The lattice is partitioned into visible
(V) and hidden (H) sectors via the checkerboard rule: site \textbf{i} is
visible if \ensuremath{\Sigma}\_k i\_k is even, hidden if odd. Every visible site has
hidden neighbors (C1). For a core-shell partition (visible = cube of
side r in a lattice of side L), timescale separation C2 and capacity
asymmetry C3 are controlled by the ratio (L \ensuremath{-} r)/r.

\begin{center}\rule{0.5\linewidth}{0.5pt}\end{center}

\hypertarget{link-1-emergent-quantum-mechanics}{%
\subsection{3. Link 1: Emergent Quantum
Mechanics}\label{link-1-emergent-quantum-mechanics}}

\textbf{Theorem 1.} \emph{The visible-sector propagator of the wave
equation with checkerboard partition satisfies G(2) \ensuremath{\neq} G(1)\ensuremath{^2} for all L \ensuremath{\geq}
4 in any dimension d \ensuremath{\geq} 1. The reduced dynamics is non-Markovian.}

\textbf{Proof.} The full-system evolution in phase space (x(t), x(t\ensuremath{-}1))
\ensuremath{\in} \ensuremath{\mathbb{R}}\^{}\{2N\} is the linear map

\[M = \begin{pmatrix} A & -I \\ I & 0 \end{pmatrix}\]

where A is the lattice adjacency matrix. The visible propagator at lag n
is G(n) = P\_V M\^{}n P\_V\^{}T. For L = 4, d = 1:

\[G(2) - G(1)^2 = \begin{pmatrix} 2 & 2 & 0 & 0 \\ 2 & 2 & 0 & 0 \\ 0 & 0 & 0 & 0 \\ 0 & 0 & 0 & 0 \end{pmatrix}\]

The entries are exactly 2, arising from the two hidden sites that each
connect the visible pair. For general L \ensuremath{\geq} 4, every pair of visible sites
at lattice distance 2 shares at least one hidden neighbor, so G(2) \ensuremath{-}
G(1)\ensuremath{^2} \ensuremath{\neq} 0. \ensuremath{\blacksquare}

\textbf{Corollary.} By the Barandes stochastic-quantum correspondence
{[}2{]}, the visible-sector dynamics admits a description as unitary
quantum mechanics with a Hermitian Hamiltonian and Born rule.

\textbf{Theorem 1b} (C2-C3 strengthening). \emph{The (i,j) entry of G(2)
\ensuremath{-} G(1)\ensuremath{^2} equals \textbar N\_H(i) \ensuremath{\cap} N\_H(j)\textbar, the number of common
hidden neighbors. For core-shell partitions, this grows monotonically as
the visible core shrinks and the hidden shell grows.} \ensuremath{\blacksquare}

\textbf{Theorem 1c} (Validity window). \emph{The emergent Hamiltonian is
valid for lags n \textless{} \ensuremath{\tau}\_B = (L \ensuremath{-} r)/(2v). The wave equation has
an exact light cone (range R = 1, no exponential tails), so information
entering the hidden shell at the core boundary takes at least \ensuremath{\tau}\_B steps
to traverse the shell and return. Before this time, the emergent
description holds; after it, returning correlations invalidate it.} \ensuremath{\blacksquare}

\textbf{Numerical confirmation.} For finite q, P-indivisibility was
confirmed via conditional mutual information against a Markov null:

\begin{longtable}[]{@{}ll@{}}
\toprule\noalign{}
System & CMI z-score \\
\midrule\noalign{}
\endhead
\bottomrule\noalign{}
\endlastfoot
1+1D, L=60, mod-256 wave eq & \textbf{31.1} \\
3+1D, L=8, mod-256 wave eq & \textbf{4.7} \\
1+1D, L=60, binary CA & \textbf{30.3} \\
3+1D, 16\ensuremath{^3}, core r=2 (511:1) & \textbf{14.9} \\
\end{longtable}

\begin{center}\rule{0.5\linewidth}{0.5pt}\end{center}

\hypertarget{link-2-emergent-spacetime}{%
\subsection{4. Link 2: Emergent
Spacetime}\label{link-2-emergent-spacetime}}

\textbf{Theorem 2.} \emph{The (site, time) partial order on a
d-dimensional lattice with L\_\ensuremath{\infty} light cone has Myrheim-Meyer dimension
converging to d + 1 as L \ensuremath{\to} \ensuremath{\infty}.}

\textbf{Proof.} The causal relation (i, t) \ensuremath{\prec} (j, t\ensuremath{^{\prime}}) holds iff t
\textless{} t\ensuremath{^{\prime}} and \textbar\textbar i \ensuremath{-} j\textbar\textbar\_\ensuremath{\infty} \ensuremath{\leq} v(t\ensuremath{^{\prime}} \ensuremath{-}
t). By the Bombelli-Lee-Meyer-Sorkin theorem {[}4{]}, the ordering
fraction converges to the value for d+1 Minkowski spacetime. Numerical
confirmation: MM dimension = 3.97 in 3+1D at L=14. \ensuremath{\blacksquare}

\begin{center}\rule{0.5\linewidth}{0.5pt}\end{center}

\hypertarget{link-3-emergent-ux210f}{%
\subsection{5. Link 3: Emergent \ensuremath{\hbar}}\label{link-3-emergent-ux210f}}

\textbf{Theorem 3} (Gap Equation). \emph{The classical horizon
temperature and the emergent quantum horizon temperature, computed
independently, must agree. This uniquely determines}

\[\hbar = \frac{c^3 \varepsilon^2}{4G}\]

\emph{connecting the lattice spacing \ensuremath{\varepsilon} to Newton's constant.
Self-consistently, \ensuremath{\varepsilon} = 2l\_P.}

\textbf{Proof.} The argument proceeds in four steps.

\emph{Step 1: Classical horizon temperature.} The lattice has a
discreteness scale \ensuremath{\varepsilon}. A local causal horizon with surface gravity \ensuremath{\kappa} has
an entropy density of one degree of freedom per minimal cell of area \ensuremath{\varepsilon}\ensuremath{^2},
so dS = dA/\ensuremath{\varepsilon}\ensuremath{^2}. The Jacobson identity {[}3{]} gives the energy flux
across the horizon as dE = (c\ensuremath{^2} \ensuremath{\kappa} / 8\ensuremath{\pi} G) dA. Applying the Clausius
relation dE = T dS:

\[T_{\text{cl}} = \frac{c^2 \varepsilon^2 \kappa}{8\pi G k_B}\]

This temperature is computed entirely within the classical substratum.
No \ensuremath{\hbar} appears.

\emph{Step 2: Boundary-only dependence.} The emergent action scale \ensuremath{\hbar} is
determined by the partition boundary, not by volumetric hidden-sector
data. This follows from spatial locality: the transition probabilities
\(T_{ij}(t)\) depend, to leading order, only on the visible-boundary
dynamics (see {[}1, §5.2{]} for the full argument). The
partition-intrinsic quantities with which to construct an action are c,
G, and \ensuremath{\varepsilon}. The unique combination with dimensions of action is \ensuremath{\hbar} = \ensuremath{\beta}
c\ensuremath{^3}\ensuremath{\varepsilon}\ensuremath{^2}/G for a dimensionless constant \ensuremath{\beta}.

\emph{Step 3: Emergent quantum temperature.} The emergent QFT of Link 1
lives on the classical background, which possesses a bifurcate Killing
horizon with surface gravity \ensuremath{\kappa}. Regularity of the Euclidean
(Wick-rotated) metric at the horizon enforces periodicity 2\ensuremath{\pi}c/\ensuremath{\kappa} in
imaginary time. Any QFT on this background --- including the
lattice-regularized one emerging from the OI framework --- must
therefore be in a KMS (thermal) state at temperature:

\[T_Q = \frac{\hbar \kappa}{2\pi c k_B}\]

with \ensuremath{\hbar} as the sole unknown. This is a theorem \emph{within} the derived
QFT, not an external import. For the lattice theory, thermal periodicity
is robust against UV modifications of the dispersion relation, with
corrections at O((\ensuremath{\varepsilon}\ensuremath{\kappa}/c)\ensuremath{^2}) {[}see references in 1, §5.2{]}.

\emph{Step 4: Thermal matching.} A thermometer at the horizon is a
visible-sector object whose click rate is computable in either
description. The classical and emergent accounts of one boundary cannot
disagree:

\[T_{\text{cl}} = T_Q \implies \frac{c^2 \varepsilon^2 \kappa}{8\pi G k_B} = \frac{\hbar \kappa}{2\pi c k_B}\]

The surface gravity \ensuremath{\kappa} cancels --- confirming that \ensuremath{\hbar} is structural
(boundary-dependent) rather than state-dependent. Solving:

\[\boxed{\hbar = \frac{c^3 \varepsilon^2}{4G}}\]

Rearranging: \ensuremath{\varepsilon}\ensuremath{^2} = 4\ensuremath{\hbar}G/c\ensuremath{^3} = 4l\_P\ensuremath{^2}, giving \ensuremath{\varepsilon} = 2l\_P. The factor 1/4 in
the Bekenstein-Hawking formula S = A/(4l\_P\ensuremath{^2}) is then derived: each cell
of area \ensuremath{\varepsilon}\ensuremath{^2} = 4l\_P\ensuremath{^2} contributes one entropy unit, so S = A/\ensuremath{\varepsilon}\ensuremath{^2} =
A/(4l\_P\ensuremath{^2}). \ensuremath{\blacksquare}

\emph{Remark.} The two temperatures are computed from independent
inputs: T\_cl from the classical substratum alone (Axioms 1--3, the
Jacobson identity, and \ensuremath{\eta} = 1/\ensuremath{\varepsilon}\ensuremath{^2}), T\_Q from the emergent QFT alone (Link
1's theorem plus Euclidean regularity). The match is non-trivial: T\_cl
could have depended on volumetric hidden-sector data (excluded by Step
2), or T\_Q could have depended on the quantum state (it does not ---
the KMS temperature is purely kinematic). That neither pathology obtains
makes the gap equation a genuine determination, not a tautology.

\begin{center}\rule{0.5\linewidth}{0.5pt}\end{center}

\hypertarget{link-4-entropy-proportional-to-area}{%
\subsection{6. Link 4: Entropy Proportional to
Area}\label{link-4-entropy-proportional-to-area}}

\textbf{Theorem 4} (Area Law). \emph{The entanglement entropy of a
spatial region V in the d-dimensional wave equation scales as S(V) = \ensuremath{\eta}
\textbar\ensuremath{\partial}V\textbar{} / \ensuremath{\varepsilon}\^{}\{d\ensuremath{-}1\}. This holds for both the linear wave
equation over \ensuremath{\mathbb{R}} and the mod-q wave equation over \ensuremath{\mathbb{Z}}/q\ensuremath{\mathbb{Z}}.}

\textbf{Proof.} \emph{Over \ensuremath{\mathbb{R}}:} The linear wave equation on a lattice is
equivalent to a system of coupled harmonic oscillators with a quadratic
Hamiltonian determined by the adjacency matrix A. This is a Gaussian
(free bosonic) system, so the area-law theorem of Eisert, Cramer, and
Plenio {[}5{]} applies directly.

\emph{Over \ensuremath{\mathbb{Z}}/q\ensuremath{\mathbb{Z}}:} The area law follows from the spatial Markov property
of nearest-neighbor dynamics. Let V° = V ~\ensuremath{\partial}V be the interior of a block
V and \ensuremath{\partial}V its boundary (sites adjacent to V\^{}c). The wave equation has
range R = 1, so x\_i(t+1) depends only on x\_j(t) for \textbar i \ensuremath{-}
j\textbar{} \ensuremath{\leq} 1 and on x\_i(t\ensuremath{-}1). For any i \ensuremath{\in} V°, all neighbors lie
inside V, so x\_i(1) is a function of initial data within V alone. With
uniform initial conditions (independent across sites), V° \ensuremath{\perp} V\^{}c
\textbar{} \ensuremath{\partial}V. By the chain rule: I(V; V\^{}c) = I(\ensuremath{\partial}V; V\^{}c) + I(V°;
V\^{}c \textbar{} \ensuremath{\partial}V) = I(\ensuremath{\partial}V; V\^{}c) \ensuremath{\leq} H(\ensuremath{\partial}V) \ensuremath{\leq} \textbar\ensuremath{\partial}V\textbar{} ·
log\ensuremath{_2} q. The mutual information scales as boundary area, not volume, for
any q \ensuremath{\geq} 2 and d \ensuremath{\geq} 1. Numerical confirmation (Appendix A.5) gives
I/\textbar\ensuremath{\partial}V\textbar{} \ensuremath{\approx} 0.3 log\ensuremath{_2} q per boundary pair, constant to
within 3\% as block width varies from 1 to 7.

Combined with the gap equation (Theorem 3), \ensuremath{\varepsilon} = 2l\_P, giving S = \ensuremath{\eta}
A/l\_P\ensuremath{^2} for d = 3. The coefficient \ensuremath{\eta} = 1/4 is fixed by the thermal
matching condition (§5, Step 4). \ensuremath{\blacksquare}

\begin{center}\rule{0.5\linewidth}{0.5pt}\end{center}

\hypertarget{link-5-lorentz-invariance}{%
\subsection{7. Link 5: Lorentz
Invariance}\label{link-5-lorentz-invariance}}

\textbf{Theorem 5} (Relativistic Dispersion). \emph{The wave equation
has exact dispersion relation cos \ensuremath{\omega} = cos k.}

\textbf{Proof.} Substituting x\_j(t) = A exp(i(kj \ensuremath{-} \ensuremath{\omega}t)):

\[e^{-i\omega} + e^{i\omega} = e^{-ik} + e^{ik} \implies \cos\omega = \cos k \qquad \square\]

For small k: \ensuremath{\omega} = \textbar k\textbar{} + O(k\ensuremath{^3}), giving relativistic
propagation with v = 1. The identity is algebraic: for the mod-q wave
equation, the same substitution using characters of \ensuremath{\mathbb{Z}}/q\ensuremath{\mathbb{Z}} yields the same
dispersion relation exactly, not as a large-q approximation.

\begin{center}\rule{0.5\linewidth}{0.5pt}\end{center}

\hypertarget{link-6-unruh-temperature}{%
\subsection{8. Link 6: Unruh
Temperature}\label{link-6-unruh-temperature}}

\textbf{Theorem 6} (Lattice Bisognano-Wichmann). \emph{For free lattice
systems whose low-energy effective theory is Lorentz-invariant, the
Bisognano-Wichmann theorem holds on the lattice: the entanglement
Hamiltonian for a half-space bipartition has the BW form (linear spatial
deformation of the physical Hamiltonian), with errors vanishing as
\ensuremath{\ell}\^{}\{\ensuremath{-}2\} {[}6, 7{]}. For free-fermion and harmonic-oscillator chains,
this has been proven analytically {[}8{]}.}

The mod-q wave equation is a coupled harmonic oscillator system ---
precisely the class for which the lattice BW theorem is proven. Its
low-energy dispersion is relativistic (Theorem 5). Therefore the reduced
state for a Rindler observer is thermal at

\[T_U = \frac{\hbar\kappa}{2\pi c k_B}\]

where \ensuremath{\kappa} is the surface gravity. The thermality of the Unruh effect is
robust against UV modifications of the dispersion relation: for the
lattice dispersion cos \ensuremath{\omega} = cos k, corrections to the Unruh temperature
scale as O(\ensuremath{\varepsilon}\ensuremath{^2}\ensuremath{\kappa}\ensuremath{^2}/c\ensuremath{^2}) {[}9, 10{]}. \ensuremath{\blacksquare}

\begin{center}\rule{0.5\linewidth}{0.5pt}\end{center}

\hypertarget{link-7-einsteins-equations}{%
\subsection{9. Link 7: Einstein's
Equations}\label{link-7-einsteins-equations}}

\textbf{Theorem 7} (Jacobson {[}3{]}). \emph{Given (i) \ensuremath{\delta}S = \ensuremath{\eta} \ensuremath{\delta}A/l\_P\ensuremath{^2}
for local causal horizons, (ii) T = \ensuremath{\hbar}\ensuremath{\kappa}/(2\ensuremath{\pi}ck\_B), and (iii) the
Raychaudhuri equation, the Clausius relation \ensuremath{\delta}Q = T\ensuremath{\delta}S yields:}

\[R_{\mu\nu} - \frac{1}{2}Rg_{\mu\nu} + \Lambda g_{\mu\nu} = \frac{8\pi G}{c^4}T_{\mu\nu}\]

The OI framework provides all three inputs: (i) from Theorem 4 --- the
area-law entropy holds for any local patch because the spatial Markov
property applies to arbitrary connected subregions; (ii) from Theorem 6
--- the Unruh temperature at each local horizon; (iii) from the emergent
(d+1)-dimensional Lorentzian manifold (Theorem 2), on which the
Raychaudhuri equation is a kinematic identity relating the expansion of
null congruences to the Ricci tensor. \ensuremath{\blacksquare}

\begin{center}\rule{0.5\linewidth}{0.5pt}\end{center}

\hypertarget{the-complete-chain}{%
\subsection{10. The Complete Chain}\label{the-complete-chain}}

\begin{longtable}[]{@{}lll@{}}
\toprule\noalign{}
Link & Statement & Proof \\
\midrule\noalign{}
\endhead
\bottomrule\noalign{}
\endlastfoot
1 & Non-Markovian \ensuremath{\to} QM & Theorems 1, 1b, 1c \\
2 & MM dimension = d+1 & Theorem 2 \\
3 & Gap equation \ensuremath{\to} \ensuremath{\hbar} & Theorem 3 \\
4 & S \ensuremath{\propto} A (area law) & Theorem 4 \\
5 & cos \ensuremath{\omega} = cos k (Lorentz invariance) & Theorem 5 \\
6 & Unruh temperature & Theorem 6 \\
7 & Einstein's equations & Theorem 7 \\
\end{longtable}

All seven links are analytical for the linear wave equation over \ensuremath{\mathbb{R}}. For
the mod-q wave equation over \ensuremath{\mathbb{Z}}/q\ensuremath{\mathbb{Z}}, all seven have analytical proofs: the
lattice BW theorem (Link 6) is proven for free-particle systems {[}8{]},
the class to which the wave equation belongs.

\begin{center}\rule{0.5\linewidth}{0.5pt}\end{center}

\hypertarget{discussion}{%
\subsection{11. Discussion}\label{discussion}}

\hypertarget{the-dynamics-selection}{%
\subsubsection{11.1 The dynamics
selection}\label{the-dynamics-selection}}

The OI axioms derive quantum mechanics \textbf{generically}: any
bijective dynamics with a coupled V/H partition produces non-Markovian
reduced dynamics (Theorem 1 for linear systems; numerical confirmation
for nonlinear CAs with z up to 30.3). Lorentz invariance is \textbf{not
generic} --- it requires wave-equation dynamics. But this choice is not
arbitrary; it is uniquely determined by three physical requirements.

\textbf{Theorem 8} (Dynamics selection). \emph{Among all second-order
reversible nearest-neighbor dynamics on a d-dimensional lattice of
alphabet size q, the requirements of (i) center independence, (ii)
spatial isotropy, and (iii) linearity uniquely select the discrete wave
equation f(neighbors) = \ensuremath{\Sigma} neighbors mod q. This holds for any q \ensuremath{\geq} 2 and
d \ensuremath{\geq} 1.}

\textbf{Proof.} The update function has the form x\_i(t+1) =
(f(neighbors of i at time t) \ensuremath{-} x\_i(t\ensuremath{-}1)) mod q. Center independence
requires that f not depend on x\_i itself: f reduces to a function of
the 2d neighboring values only. Spatial isotropy requires f to be
symmetric under all permutations of the neighbors that correspond to
lattice symmetries. Linearity over \ensuremath{\mathbb{Z}}/q\ensuremath{\mathbb{Z}} requires f(a + a') = f(a) +
f(a') for neighbor configurations a, a'. The unique function satisfying
all three is f = \ensuremath{\alpha}(x\ensuremath{_1} + x\ensuremath{_2} + \ldots{} + x\_\{2d\}) mod q for some
constant \ensuremath{\alpha}. The propagation speed is v = \ensuremath{\alpha}, so \ensuremath{\alpha} = 1 gives the maximum
lattice-scale speed. This is the discrete wave equation. \ensuremath{\blacksquare}

Each requirement has a physical justification:

\textbf{Center independence is necessary for emergent QM.} Center
dependence suppresses P-indivisibility for linear rules over \ensuremath{\mathbb{Z}}/q\ensuremath{\mathbb{Z}}: Rule
150 (f = l \ensuremath{\oplus} c \ensuremath{\oplus} r) has zero P-indivisibility (CMI z = 0.0) vs Rule 90's
z = 7.2, despite identical propagation speed. The mechanism ---
information screening by the visible self-coupling --- is proven
analytically in §11.2 and confirmed numerically through q = 64 (Appendix
A.4).

\textbf{Spatial isotropy is required for Lorentz invariance.} Without
left-right (or rotational) symmetry, the dynamics has a preferred
spatial direction, breaking the isotropy that Lorentz invariance
demands.

\textbf{Linearity gives maximum propagation speed.} Among all
center-independent, isotropic dynamics, the linear wave equation
achieves the maximum propagation speed v = 1 (the lattice speed of
light). The dispersion relation cos \ensuremath{\omega} = cos k (Theorem 5) gives exact
group velocity v\_g = sin k / sin \ensuremath{\omega} \ensuremath{\to} 1 as k \ensuremath{\to} 0. Nonlinear
center-independent dynamics generically have v \textless{} 1, confirmed
by the exhaustive scan of all 256 binary rules: the 53 QM-producing
rules have mean speed 0.33, while only Rule 90 achieves v = 1
(anti-correlation r = \ensuremath{-}0.60, p \textless{} 10\ensuremath{^{-}}\ensuremath{^4}).

The derivation chain for GR therefore requires wave-equation dynamics,
uniquely selected by center independence, isotropy, and linearity.
Linearity is itself selected by three independent criteria (§11.2):
maximum propagation speed, maximum P-indivisibility, and horizon
thermality.

\hypertarget{caveats}{%
\subsubsection{11.2 Caveats}\label{caveats}}

\textbf{The continuum limit and Jacobson stability.} The derivation
chain works on a finite lattice but invokes continuum results in the
final link (Jacobson's thermodynamic derivation). The concern is whether
lattice corrections to Lorentz invariance invalidate the result. This
can be resolved by decomposing Jacobson's derivation into its four
inputs and checking each independently:

\emph{(i) \ensuremath{\delta}S = \ensuremath{\delta}A/\ensuremath{\varepsilon}\ensuremath{^2} (area-law entropy).} Proven for Gaussian systems
over \ensuremath{\mathbb{R}} {[}5{]} and for any nearest-neighbor dynamics over \ensuremath{\mathbb{Z}}/q\ensuremath{\mathbb{Z}} via the
spatial Markov property (Theorem 4). No Lorentz invariance enters.
\textbf{Exact on the lattice.}

\emph{(ii) T = \ensuremath{\hbar}\ensuremath{\kappa}/(2\ensuremath{\pi}ck\_B) (Unruh temperature).} This requires the
Bisognano-Wichmann theorem. On the lattice, the BW form of the
entanglement Hamiltonian is proven analytically for free-particle
systems {[}8{]}, confirmed numerically across multiple universality
classes {[}6, 7{]}, and shown to be robust against UV modifications of
the dispersion relation with corrections at O(\ensuremath{\varepsilon}\ensuremath{^2}\ensuremath{\kappa}\ensuremath{^2}/c\ensuremath{^2}) {[}9, 10{]}.
\textbf{Approximate: O(\ensuremath{\varepsilon}\ensuremath{^2}\ensuremath{\kappa}\ensuremath{^2}/c\ensuremath{^2}) corrections.}

\emph{(iii) \ensuremath{\delta}Q = \ensuremath{\int} T\_\ensuremath{\mu}\ensuremath{\nu} k\^{}\ensuremath{\mu} d\ensuremath{\Sigma}\^{}\ensuremath{\nu} (energy flux).} This is a
definition on the emergent manifold (Theorem 2). The stress-energy
tensor is constructed from the emergent fields; the integral is over a
null generator of the local causal horizon. No Lorentz invariance is
required. \textbf{Exact on the emergent manifold.}

\emph{(iv) Raychaudhuri equation.} The identity d\ensuremath{\theta}/d\ensuremath{\lambda} = \ensuremath{-}(1/2)\ensuremath{\theta}\ensuremath{^2} \ensuremath{-}
\ensuremath{\sigma}\_\{ab\}\ensuremath{\sigma}\^{}\{ab\} \ensuremath{-} R\_\ensuremath{\mu}\ensuremath{\nu} k\^{}\ensuremath{\mu} k\^{}\ensuremath{\nu} is kinematic: it holds on any
pseudo-Riemannian manifold as a consequence of the definition of
curvature. No Lorentz invariance is required. \textbf{Exact on the
emergent manifold.}

Since inputs (i), (iii), and (iv) are lattice-exact, corrections to
Einstein's equations enter only through input (ii):

\[G_{\mu\nu} + \Lambda g_{\mu\nu} = \frac{8\pi G}{c^4} T_{\mu\nu} \times \left[1 + \mathcal{O}\!\left(\frac{\varepsilon^2 \kappa^2}{c^2}\right)\right]\]

For the cosmological horizon, \ensuremath{\varepsilon}\ensuremath{\kappa}/c = 2l\_P · H/c \textasciitilde{}
10\ensuremath{^{-}}\ensuremath{^6}\ensuremath{^1}, giving corrections of order 10\ensuremath{^{-}}\ensuremath{^1}\ensuremath{^2}\ensuremath{^2}. For solar-mass black holes,
\ensuremath{\varepsilon}\ensuremath{\kappa}/c \textasciitilde{} 10\ensuremath{^{-}}\ensuremath{^3}\ensuremath{^8}, giving corrections of order 10\ensuremath{^{-}}\ensuremath{^7}\ensuremath{^6}. The
corrections are negligible for any horizon much larger than the Planck
scale.

\textbf{Center independence and emergent QM.} Theorem 8 relies on center
independence as necessary for emergent QM. Numerical experiments
(Appendix A.4) confirm this for every q from 2 to 64. The following
argument proves the mechanism on the full lattice, not just a single
site.

For the checkerboard partition, the adjacency matrix A is bipartite:
A\_VV = 0 (no visible--visible edges). A center-dependent rule replaces
A with A + I, giving (A + I)\_VV = I.

\emph{CI case (A\_VV = 0):} For each visible site i, the update
x\_i(t+1) = {[}h\_\{i\ensuremath{-}1\}(t) + h\_\{i+1\}(t) \ensuremath{-} x\_i(t\ensuremath{-}1){]} mod q does
not contain x\_i(t). Therefore X\_V(t+1) is a function of (H(t),
X\_V(t\ensuremath{-}1)) that carries no information about X\_V(t). Conditioning on
X\_V(t+1) constrains H(t) but does not screen X\_V(t) from the future.
Since H(t+1) depends on X\_V(t), the path X\_V(t) \ensuremath{\to} H(t+1) \ensuremath{\to} X\_V(t+2)
bypasses the present: the process is non-Markovian.

*CD case ((A+I)\emph{VV = I):* The update x\_i(t+1) = {[}x\_i(t) +
h\_\{i\ensuremath{-}1\}(t) + h\_\{i+1\}(t) \ensuremath{-} x\_i(t\ensuremath{-}1){]} mod q contains x\_i(t)
explicitly. Conditioning on (X\_V(t), X\_V(t+1)) gives, for each visible
site i, the constraint h}\{i\ensuremath{-}1\}(t) + h\_\{i+1\}(t) = x\_i(t+1) \ensuremath{-}
x\_i(t) + x\_i(t\ensuremath{-}1) mod q. This is a system of \textbar V\textbar{}
linear equations in \textbar H\textbar{} unknowns over \ensuremath{\mathbb{Z}}/q\ensuremath{\mathbb{Z}}, with
coefficient matrix A\_VH. When this system has a unique solution (which
it does generically for q prime), the hidden state H(t) is fully
determined by (X\_V(t), X\_V(t+1), X\_V(t\ensuremath{-}1)). Since the full state
determines all future evolution, the augmented process Y\_V(t) =
(X\_V(t), X\_V(t\ensuremath{-}1)) is Markov. Numerical verification: the
augmented-state CMI for CD is suppressed by 10--30\ensuremath{\times} relative to CI at
all tested q from 3 to 7.

The mechanism is information screening: center coupling makes X\_V(t+1)
an explicit function of X\_V(t), allowing the present to determine the
hidden state and screen the past from the future. Without center
coupling (A\_VV = 0), the present carries no information about the
visible past, so the hidden sector retains unscreened correlations. Over
\ensuremath{\mathbb{R}}, the screening does not operate --- G(2) \ensuremath{-} G(1)\ensuremath{^2} is identical for CI
and CD --- because the mod operation is essential to the conditional
entropy structure.

\textbf{Linearity selection.} The linearity requirement in Theorem 8 is
selected by three independent criteria, none imposed as an axiom. First,
over \ensuremath{\mathbb{R}}, nonlinear CI dynamics have propagation speed v \textless{} 1, so
linearity gives maximum speed (§11.1). Second, over \ensuremath{\mathbb{Z}}/q\ensuremath{\mathbb{Z}}, among all
linear CI dynamics f = \ensuremath{\alpha}(l + r) mod q with prime q \ensuremath{\geq} 5, \ensuremath{\alpha} = 1 uniquely
maximizes P-indivisibility (Appendix A.6): z-scores of 3--11 vs z \ensuremath{\approx} 0
for all other \ensuremath{\alpha}. Third, modified (nonlinear) dispersion relations break
the thermality of the Unruh effect {[}9, 10{]} --- the KMS condition
holds only for linear dispersion \ensuremath{\omega} = \textbar k\textbar. Since
Jacobson's derivation requires thermal equilibrium at horizons, only the
wave equation supports the full chain. All three criteria converge on
the same dynamics.

\textbf{The \ensuremath{\eta} = 1/4 coefficient.} The Bekenstein-Hawking formula S =
A/(4l\_P\ensuremath{^2}) follows from \ensuremath{\varepsilon} = 2l\_P (the gap equation): one entropy unit
per cell of area \ensuremath{\varepsilon}\ensuremath{^2} = 4l\_P\ensuremath{^2}. This cannot be confirmed via entanglement
entropy because of the species problem --- the entanglement entropy
coefficient is species-dependent (Srednicki found \ensuremath{\eta}\_ent \ensuremath{\approx} 0.295 per
scalar field), while the BH entropy is not. The thermal matching route
avoids this by counting classical boundary DOF (one per \ensuremath{\varepsilon}\ensuremath{^2} cell,
independent of field content) rather than any particular field's vacuum
entanglement.

\textbf{Jacobson's local equilibrium.} Required only for input (ii)
above. For the wave equation, the lattice BW theorem {[}8{]} proves the
entanglement Hamiltonian has the BW form exactly, establishing local
thermal equilibrium at each causal horizon up to O(\ensuremath{\varepsilon}\ensuremath{^2}\ensuremath{\kappa}\ensuremath{^2}/c\ensuremath{^2}) corrections.

\textbf{Field content.} The wave equation is a scalar (bosonic) system.
The Standard Model contains fermions and gauge fields, which are absent
from the lattice construction. This does not affect the GR derivation
chain: the Jacobson route requires only thermodynamic inputs (entropy
proportional to area, the Unruh temperature, and the Raychaudhuri
equation), none of which depend on the emergent field content. A
companion paper {[}11{]} shows that fermionic and gauge degrees of
freedom arise from the same lattice substrate: the wave equation factors
into staggered Dirac operators (Susskind factorization), center
independence enforces chiral symmetry, and the multi-component extension
with \(K = 6\) internal components produces the Standard Model gauge
group \(\mathrm{SU}(3) \times \mathrm{SU}(2) \times \mathrm{U}(1)\) from
the cubic group decomposition of the link directions.

\begin{center}\rule{0.5\linewidth}{0.5pt}\end{center}

\hypertarget{conclusion}{%
\subsection{12. Conclusion}\label{conclusion}}

The OI framework {[}1{]} proves that QM emerges from embedded
observation. This paper tested the framework's rigidity by constructing
an explicit lattice realization and checking whether the dynamics
selected by the QM requirement also produces the inputs for general
relativity.

The dynamics selection is unique: among all second-order reversible
nearest-neighbor rules, center independence (necessary for emergent QM),
spatial isotropy, and linearity select the discrete wave equation alone
(Theorem 8). Linearity is itself selected by three independent criteria
--- maximum propagation speed, maximum P-indivisibility, and horizon
thermality --- leaving no free parameters.

This dynamics then passes seven independent checks (Links 1--7), all
analytically proven for both \ensuremath{\mathbb{R}} and \ensuremath{\mathbb{Z}}/q\ensuremath{\mathbb{Z}}: non-Markovian reduced dynamics,
correct emergent spacetime dimension, the gap equation for \ensuremath{\hbar}, area-law
entropy, Lorentz-invariant dispersion, the Unruh temperature, and all
inputs to Jacobson's derivation of Einstein's equations. The continuum
limit is controlled: three of four Jacobson inputs are lattice-exact,
and the fourth carries corrections of O(\ensuremath{\varepsilon}\ensuremath{^2}\ensuremath{\kappa}\ensuremath{^2}/c\ensuremath{^2}) --- negligible for all
physical horizons (§11.2).

Each link is a point where the framework could have failed and did not.
The structural consistency --- one dynamics, selected by one criterion,
passing seven tests with no tuning --- constitutes independent
corroborative evidence for the OI framework's central claim.

\begin{center}\rule{0.5\linewidth}{0.5pt}\end{center}

\hypertarget{acknowledgements}{%
\subsection{Acknowledgements}\label{acknowledgements}}

During the preparation of this work, the author used Claude Opus 4.6
(Anthropic) to assist in drafting, numerical simulations, literature
review, and refining argumentation. The author reviewed and edited all
content and takes full responsibility for the publication.

\begin{center}\rule{0.5\linewidth}{0.5pt}\end{center}

\hypertarget{references}{%
\subsection{References}\label{references}}

{[}1{]} A. Maybaum, ``The Incompleteness of Observation,'' submitted to
Foundations of Physics (2026).

{[}2{]} J. Barandes, ``The stochastic-quantum correspondence,''
arXiv:2302.10778 (2023).

{[}3{]} T. Jacobson, ``Thermodynamics of Spacetime: The Einstein
Equation of State,'' Phys. Rev.~Lett. 75, 1260 (1995).

{[}4{]} L. Bombelli, J. Lee, D. Meyer, R. Sorkin, ``Space-time as a
causal set,'' Phys. Rev.~Lett. 59, 521 (1987).

{[}5{]} J. Eisert, M. Cramer, M.B. Plenio, ``Area laws for the
entanglement entropy,'' Rev.~Mod. Phys. 82, 277 (2010).

{[}6{]} G. Giudici, T. Mendes-Santos, P. Calabrese, M. Dalmonte,
``Entanglement Hamiltonians of lattice models via the Bisognano-Wichmann
theorem,'' Phys. Rev.~B 98, 134403 (2018).

{[}7{]} J. Zhang, P. Calabrese, M. Dalmonte, M.A.~Rajabpour, ``Lattice
Bisognano-Wichmann modular Hamiltonian in critical quantum spin
chains,'' SciPost Phys. Core 2, 007 (2020).

{[}8{]} V. Eisler, ``On the Bisognano-Wichmann entanglement Hamiltonian
of nonrelativistic fermions,'' J. Stat. Mech. (2025) 013101;
arXiv:2410.16433 (2024).

{[}9{]} W.G. Unruh, ``Sonic analogue of black holes and the effects of
high frequencies on black hole evaporation,'' Phys. Rev.~D 51, 2827
(1995).

{[}10{]} L.J. Garay, J.R. Anglin, J.I. Cirac, P. Zoller, ``Influence of
the dispersion relation on the Unruh effect,'' Phys. Rev.~D 103, 085010
(2021).

{[}11{]} A. Maybaum, ``Why These Particles: Standard Model Structure
from Observational Incompleteness,'' companion paper (2026).

\begin{center}\rule{0.5\linewidth}{0.5pt}\end{center}

\hypertarget{appendix-numerical-confirmation-of-dynamics-selection}{%
\subsection{Appendix: Numerical Confirmation of Dynamics
Selection}\label{appendix-numerical-confirmation-of-dynamics-selection}}

Theorem 8 is an algebraic result. The following numerical data, from an
exhaustive scan of all 256 second-order reversible binary CA rules
(L=50, T=80), confirms the key claims empirically.

\hypertarget{a.1-selection-counts}{%
\subsubsection{A.1 Selection counts}\label{a.1-selection-counts}}

\begin{longtable}[]{@{}ll@{}}
\toprule\noalign{}
Criterion & Rules passing (of 256) \\
\midrule\noalign{}
\endhead
\bottomrule\noalign{}
\endlastfoot
Nontrivial dynamics & 254 \\
Robust QM (min z \textgreater{} 2, multi-seed) & 53 \\
Ballistic propagation (speed \textgreater{} 0.7, R\ensuremath{^2} \textgreater{} 0.9)
& 24 \\
Area-law entropy scaling & 18 \\
QM \ensuremath{\cap} ballistic & 1--2 \\
\end{longtable}

\hypertarget{a.2-algebraic-characterization}{%
\subsubsection{A.2 Algebraic
characterization}\label{a.2-algebraic-characterization}}

\begin{longtable}[]{@{}lll@{}}
\toprule\noalign{}
Property & Count (of 256) & Physical interpretation \\
\midrule\noalign{}
\endhead
\bottomrule\noalign{}
\endlastfoot
Center-independent & 16 & Hidden sector mediates V-V interactions \\
Left-right symmetric & 64 & Spatial isotropy \\
Linear over GF(2) & 8 & Gaussian structure \\
All three (excluding trivial) & 1 (Rule 90) & Wave equation \\
\end{longtable}

\hypertarget{a.3-rule-90-vs-rule-150}{%
\subsubsection{A.3 Rule 90 vs Rule 150}\label{a.3-rule-90-vs-rule-150}}

\begin{longtable}[]{@{}lll@{}}
\toprule\noalign{}
Property & Rule 90: f = l \ensuremath{\oplus} r & Rule 150: f = l \ensuremath{\oplus} c \ensuremath{\oplus} r \\
\midrule\noalign{}
\endhead
\bottomrule\noalign{}
\endlastfoot
Propagation speed & 1.000 & 1.000 \\
CMI z (L=50) & 7.2 \ensuremath{\pm} 1.1 & 0.0 \ensuremath{\pm} 0.6 \\
Center-independent & Yes & No \\
\end{longtable}

The sole difference --- center dependence --- determines whether
emergent QM exists.

\hypertarget{a.4-center-independence-beyond-q-2}{%
\subsubsection{A.4 Center independence beyond q =
2}\label{a.4-center-independence-beyond-q-2}}

The following data extends the Rule 90/150 comparison to arbitrary
alphabet size q. For each q, we compare the linear wave equation f = l +
r mod q (center-independent, CI) against its center-dependent analog f =
l + c + r mod q (CD). The non-Markovianity statistic is the z-score of
\textbar\textbar G(2) \ensuremath{-} G(1)\ensuremath{^2}\textbar\textbar\_F against a time-shuffled
null (L = 60, T = 1000, 30 seeds).

\begin{longtable}[]{@{}llll@{}}
\toprule\noalign{}
q & CI z-score & CD z-score & Both v = 1 \\
\midrule\noalign{}
\endhead
\bottomrule\noalign{}
\endlastfoot
2 & 4.1 & \ensuremath{-}0.1 & Yes \\
3 & 6.8 & \ensuremath{-}0.1 & Yes \\
4 & 10.9 & 0.1 & Yes \\
5 & 8.6 & 0.1 & Yes \\
7 & 8.6 & 0.0 & Yes \\
8 & 16.4 & 0.0 & Yes \\
11 & 7.4 & 0.0 & Yes \\
13 & 5.6 & 0.1 & Yes \\
16 & 16.9 & 0.0 & Yes \\
32 & 8.4 & 0.0 & Yes \\
64 & 3.1 & 0.0 & Yes \\
\end{longtable}

For every q from 2 to 64, the center-independent wave equation produces
significant non-Markovianity while the center-dependent analog does not,
despite identical propagation speed. Among random (nonlinear) rules at q
= 3, both CI and CD rules produce P-indivisibility (47/50 and 40/50 with
z \textgreater{} 2), consistent with the general characterization
theorem {[}1{]}. Center independence is necessary specifically for the
linear dynamics.

\hypertarget{a.5-area-law-for-the-mod-q-wave-equation-in-2d}{%
\subsubsection{A.5 Area law for the mod-q wave equation in
2D}\label{a.5-area-law-for-the-mod-q-wave-equation-in-2d}}

The area law (Theorem 4) is proven analytically for both \ensuremath{\mathbb{R}} (Gaussian
argument {[}5{]}) and \ensuremath{\mathbb{Z}}/q\ensuremath{\mathbb{Z}} (spatial Markov property). The following data
provides independent numerical confirmation at finite q. On a 16 \ensuremath{\times} 16
periodic lattice with T = 8000 timesteps averaged over 5 initial
conditions, we measure the per-boundary-pair mutual information across
the right edge of a strip of width w:

\begin{longtable}[]{@{}lllllllll@{}}
\toprule\noalign{}
q & w = 1 & w = 2 & w = 3 & w = 4 & w = 5 & w = 6 & w = 7 & slope \\
\midrule\noalign{}
\endhead
\bottomrule\noalign{}
\endlastfoot
4 & 0.271 & 0.287 & 0.271 & 0.304 & 0.261 & 0.275 & 0.263 & \ensuremath{-}0.002 \\
8 & 0.714 & 0.727 & 0.712 & 0.723 & 0.701 & 0.711 & 0.692 & \ensuremath{-}0.004 \\
16 & 1.352 & 1.348 & 1.329 & 1.334 & 1.340 & 1.341 & 1.336 & \ensuremath{-}0.002 \\
32 & 2.118 & 2.110 & 2.095 & 2.094 & 2.101 & 2.097 & 2.103 & \ensuremath{-}0.002 \\
64 & 2.963 & 2.955 & 2.946 & 2.943 & 2.951 & 2.951 & 2.956 & \ensuremath{-}0.001 \\
\end{longtable}

Under an area law, the per-pair MI is independent of strip width
(constant rows). Under a volume law, it would grow proportionally to w
(a 600\% increase from w = 1 to w = 7). The observed variation is less
than 3\% at every q, with slopes consistent with zero. The area law
holds for the mod-q wave equation at all tested alphabet sizes.

\hypertarget{a.6-ux3b1-selection-wave-equation-maximizes-p-indivisibility}{%
\subsubsection{A.6 \ensuremath{\alpha}-selection: wave equation maximizes
P-indivisibility}\label{a.6-ux3b1-selection-wave-equation-maximizes-p-indivisibility}}

Among all linear center-independent dynamics f = \ensuremath{\alpha}(l + r) mod q, the
wave equation (\ensuremath{\alpha} = 1, and its equivalent \ensuremath{\alpha} = q\ensuremath{-}1) uniquely maximizes
P-indivisibility for prime q. The statistic is the z-score of
\textbar\textbar G(2) \ensuremath{-} G(1)\ensuremath{^2}\textbar\textbar\_F against a time-shuffled
null (L = 60, T = 800, 15 seeds). For each q, the table shows \ensuremath{\alpha} = 1
vs.~the best-performing alternative \ensuremath{\alpha}.

\begin{longtable}[]{@{}lllll@{}}
\toprule\noalign{}
q & \ensuremath{\alpha} = 1 z-score & Best other \ensuremath{\alpha} & Best other z & \ensuremath{\alpha} = 1 wins? \\
\midrule\noalign{}
\endhead
\bottomrule\noalign{}
\endlastfoot
5 & 7.1 & 2 & 0.9 & Yes \\
7 & 4.2 & 5 & 0.4 & Yes \\
11 & 7.7 & 5 & 0.2 & Yes \\
13 & 2.8 & 10 & 0.1 & Yes \\
17 & 1.7 & 14 & 0.5 & Yes \\
19 & 2.4 & 13 & 0.2 & Yes \\
23 & 5.8 & 13 & 0.1 & Yes \\
\end{longtable}

For every prime q tested, \ensuremath{\alpha} = 1 is the unique maximizer of
P-indivisibility (excluding the equivalent \ensuremath{\alpha} = q\ensuremath{-}1). For composite q (4,
8, 16), some non-unit \ensuremath{\alpha} values match \ensuremath{\alpha} = 1 by effectively reducing the
system to a smaller prime-power alphabet. The wave equation is selected
by maximum P-indivisibility, providing an independent criterion that
converges with the maximum-speed argument over \ensuremath{\mathbb{R}}.

\end{document}
